\section{Operational refinements of the retained transport results}\label{sec:retained}
This appendix records precise supplements to the complete relative-report article in Companion G. Its exact circle and projective-orbit arguments are not shortened or replaced. The refinements do not extend a state-independent theorem to informative reports; that extension is proved separately in the main text.

\subsection{Scheduled and autonomous information patterns}
\begin{center}\small
\begin{tabular}{p{.22\textwidth}p{.32\textwidth}p{.33\textwidth}}
\toprule
 & Scheduled observer & Autonomous observer\\
\midrule
At update & Old label, current report, fresh randomness and phase & Old label, current report and fresh randomness only\\[3pt]
Rows & A prescribed row at each phase & One reused row; phase must be encoded in labels\\[3pt]
Budget & $|S_t|\le m_t$ at each phase & Initial, intermediate and halt labels all count in $M$\\[3pt]
Deadline & Externally scheduled & Must halt exactly internally; reports remain available\\[3pt]
Relative circle & Excess $1-\prod_t\beta_{m_t}$ & $1+\sum_tm_t\le M$, balanced widths attain optimum\\[3pt]
Stable Bayesian class & Bounds~\eqref{eq:curvature-recursion} or~\eqref{eq:pathwise-recursion} & Excess comparable to $(M-n)^{-2/d}$\\
\bottomrule
\end{tabular}
\end{center}
On the circle, $\beta_m=(\sin(\pi/m)/(\pi/m))^2$ with $\beta_1=0$. For $M-1=nq+r$, $0\le r<n$, the retained autonomous excess is $1-\beta_q^{n-r}\beta_{q+1}^r$, with feasibility $M\ge n+1$. Exact real comparisons to regular arc boundaries are Borel operations. A bit program must use the separately charged report-precision oracle and approximate constants; it is not an exact implementation of arbitrary real arithmetic.

The static fractional-cell rearrangement used there follows without an assumption of deterministic cells. For a measurable membership $0\le f\le1$ of circle mass $p$, rotate the first moment to the real axis. Among such memberships, the integral of $f(\theta)\cos\theta$ is maximized by taking the largest values of $\cos\theta$, namely the centered arc of mass $p$. Subtract its indicator: the difference has zero integral and opposite sign to $\cos\theta-\cos(\pi p)$, proving the inequality directly. The magnitude of the normalized first moment is at most $\sin(\pi p)/\pi$. This is the rearrangement step in the retained static proof, consistent with the circle quantization literature~\cite{circle}.

Its exact local saturation statement requires an attaining direction, an attaining row quantizer, and identical centroids at every common outgoing label. Static compact quantization attains its center infimum, but a directional coefficient defined over a larger family may only be an infimum. Exact circle arcs supply global compatible recursion; the higher-dimensional orbit cover supplies only a matching rate. A zero-energy row contributes no constraint, and a later zero product factor does not force earlier equalities. Repeated or unused codepoints give at most the declared number of labels. The response set is the closed unit ball, or the appropriate convex hull, so the scaled terminal readout is legal.

\subsection{A modular spectral projection estimate}
\begin{lemma}[Uniform gap and projector control]\label{lem:orbit-gap}
Let $V$ be a traceless Hermitian $D\times D$ matrix with $\|V\|_{\rm HS}=1$, and write its eigenvalues decreasingly. There is an index $j(V)\in\{1,\ldots,D-1\}$ with
\[
 \lambda_{j(V)}-\lambda_{j(V)+1}\ge g_D:=\frac1{\sqrt D(D-1)}.
\]
For any two conjugates $W,W'$ of $V$, and their top-$j(V)$ projections $P,P'$, one has
\begin{equation}\label{eq:projector-bound}
 \|P-P'\|_{\rm HS}\le g_D^{-1}\|W-W'\|_{\rm HS}.
\end{equation}
Consequently there exist $C_D,r_D>0$ such that for every such $V$, every Hermitian center $Z$, and $0<r\le r_D$,
\[
 h\{U:\|UVU^*-Z\|_{\rm HS}\le r\}\le C_D r^{2(D-1)},
\]
where $h$ is Haar probability. The constants are uniform over repeated eigenvalues.
\end{lemma}
\begin{proof}
One eigenvalue has absolute value at least $D^{-1/2}$ and trace zero places another on the other side of zero. Their range is at least $D^{-1/2}$, so an adjacent gap is at least $g_D$. Choose the smallest index meeting this bound to define $j(V)$ measurably. The top projection is unique despite possible repeated eigenvalues within either separated cluster.

The spectral gap implies
\[
 \operatorname{tr}(W(P-P'))\ge g_D(j-\operatorname{tr}(PP')),
 \quad
 \operatorname{tr}(W'(P'-P))\ge g_D(j-\operatorname{tr}(PP')).
\]
For the first inequality, expand in an eigenbasis of $W$; the trace moved out of its top subspace equals the trace moved in below it, and each such unit loses at least $g_D$. Adding the two inequalities and using
$\|P-P'\|_{\rm HS}^2=2(j-\operatorname{tr}(PP'))$ proves~\eqref{eq:projector-bound} by Cauchy--Schwarz, including its zero case.

If the radius-$r$ ball meets the orbit, fix one conjugate in it. All projections of other conjugates in the ball lie within $2r/g_D$ of its top projection. Haar top projections are uniform on the complex Grassmannian of rank $j$. In its graph chart $A\in\mathbb C^{(D-j)\times j}$, normalizing a complex Gaussian frame and writing its lower block as $AB$ gives chart density proportional to $\det(I+A^*A)^{-D}$. To verify the exponent, the substitution in the lower block has real Jacobian $|\det B|^{2(D-j)}$; integrating over the upper block after the change $(I+A^*A)^{1/2}B$ contributes determinant power $-D$. Near the origin the density is bounded and projector distance is comparable to $\|A\|_{\rm HS}$. A radius-$s$ neighborhood therefore has mass at most $C_{D,j}s^{2j(D-j)}$. Unitary invariance centers the same estimate anywhere. Since $2j(D-j)\ge2(D-1)$, and there are finitely many $j$, take the maximum of the constants and the minimum radius. This proves the uniform claim without assuming a simple spectrum.
\end{proof}
For a rank-one projector, a Borel section modulo phase is explicit: choose the first index $j$ with $P_{jj}\ge1/D$ and set $v=P e_j/\sqrt{P_{jj}}$. Then $v_j$ is real positive and $vv^*=P$. A finite-chart Gram--Schmidt rule, choosing the first nonzero residual at each step, completes $v$ to a unitary; adjusting the last unused column phase gives determinant one when needed. The charts cover the orbit and their tie sets are Borel. The projective dimension in Companion G is $2(D-1)$ for the declared linear audit, not $D^2-1$ by fiat. Its observer is classical and exact higher-dimensional product saturation is not asserted.

For the contrast normalization used there, $\rho=I_D/D+r c_D X$ with $c_D=\sqrt{1-1/D}$ and $k_q=D^2-1$. Selecting one Hilbert--Schmidt basis effect after commitment gives
\[
 \E[S e_I\mid\rho]=(r c_D/k_q)x.
\]
The basis operator norm need only be at most one. Constants depending on $D$ are distinct from those depending on $r$ (or other physical contrast); the scored audit bound deteriorates as contrast tends to zero. Sphere coordinate dimension is denoted $k_s=d_s+1$, not the same symbol as this quantum audit dimension.

\subsection{Singular mass, temporal access, and digital conventions}
For retained hidden Haar reports the actual predictive law is $w_n\mu+(1-w_n)\delta_0$. Its full-history baseline is $B_n=V-w_n$; total risk is that baseline plus the surviving-mass compression loss, not the un-erased baseline $V-1$ plus a renormalized law. Inactive labels and their phase copies are counted. Lemma~\ref{lem:clock} supplies the common positive-path-support version of the deadline argument when phase report probabilities differ but their null sets agree. A visible suppression pattern with disjoint supports can reveal a clock and is outside that lower bound.

The retained fixed-start access theorem can be sharpened for one fixed model $x$ by applying its family theorem to $\{x\}$; the remainder is then the singleton remainder $B_N(\{x\})$, with its own constants, not the original family remainder silently relabeled. The initialization calls, Boolean gate and prefix clock are extra resources of that tester and remain outside the original observer budget. An ungated scalar discrepancy can still cancel even when all rows are accessible. The corrected physical crossing convention is $C_N\le N$ and the older shifted index is $K_N=C_N+1\le N+1$; the predictable starts are $T_{j-1}<N$.

For finite-bit circle execution the mismatch bound is $\min\{1,C(M-1)h\}$, not its uncapped linear expression. Sine, cosine and response constants can be shared read-only data; copying them into every row is a different program count. Complete-law simulator comparison includes the target's fresh randomization. Protected-audit lower bounds based on a $[-1,1]$ statistic use the factor two under the stated TV convention. These are finite-horizon results without a regenerative decomposition, not assertions that a physical process has no preparation. The corrected temporal companion is distinguished from its originally reviewed source; both original source and report remain preserved independently.
